%% ma: L12-B08-0118
%% lop: 12
%% chuong: 2
%% bai: 8
%% dang: 12.8.8
%% loai: TN4
%% muc: TH
%% dap_an: "B"
%% nguon: "CĐ2.3.pdf giáo viên gửi 27/10/2026 (M07, Lê Minh Khôi – Nguyễn Anh Khôi, Chương II Chủ đề 3 Biểu thức tọa độ của các phép toán vectơ), Câu 30"
%% kien_thuc: "Bài 8 (biểu thức tọa độ của phép toán vectơ)"
%% trang_thai: chinh-thuc
\begin{ex}
Trong không gian $Oxyz$, cho các điểm $A(2;0;-1)$, $B(0;3;-2)$, $C(-4;-5;0)$. Tọa độ của điểm $M$ thỏa mãn $\overrightarrow{MA}+2\overrightarrow{MB}+3\overrightarrow{MC}=\vec 0$ là
\choice{$M\left(\dfrac53;-\dfrac32;\dfrac56\right)$}{\True $M\left(-\dfrac53;-\dfrac32;-\dfrac56\right)$}{$M\left(-\dfrac53;-\dfrac32;\dfrac56\right)$}{$M\left(\dfrac53;-\dfrac32;-\dfrac56\right)$}
\loigiai{$M=\dfrac{A+2B+3C}6=\left(\dfrac{-10}6;\dfrac{-9}6;\dfrac{-5}6\right)=\left(-\dfrac53;-\dfrac32;-\dfrac56\right)$.}
\end{ex}
